% SGA 6, Exposición VI: fragmento español alineado con la fuente.
% Alcance: sección 2, número 2.3.2, líneas 8945--9013 inclusive.
% Fuente francesa: sga6_fr_r1/src/fr.tex.
% SHA-256 de la fuente francesa completa:
% 0D2158B9FE3B466309A92165727A0B191422E05FBAA0DACD5806E8322C54BA1A.
\providecommand{\sisi}[2]{#2}

\textbf{2.3.2.} Sean ahora \(B\) una \(\mathcal O_S\)-álgebra graduada cuasicoherente, \(M\) un \(B\)-módulo graduado y \(X=\Proj(B)\). Para todo abierto afín \(U\) de \(S\), y toda sección \(s\) de \(B\) sobre \(U\), homogénea de grado \(d>0\), denotemos por \(M_{(s),U}\) la componente homogénea de grado \(0\) del módulo de fracciones \(M_s\); se tiene, pues:
\[
  M_{(s),U}=
  \bigl((M|U)\otimes_{(B|U)}(B|U)_s\bigr)_0 .
\]
Si \(V\) es un abierto afín de \(U\), entonces \(M_{(s),V}\) es la restricción a \(V\) de \(M_{(s),U}\).

Consideremos la base de abiertos de \(X\) formada por los abiertos afines:
\[
  D_+(U,s)\simeq \Spec\bigl((B|U)_{(s)}\bigr) .
\]
Denotemos por \(f\) el morfismo estructural \(X\to S\), y por \(f_{U,s}\) el morfismo afín:
\[
  f_{U,s}:D_+(U,s)\longrightarrow U .
\]
Pongamos entonces:
\[
  \Proj_0(M)_{U,s}
  =
  \bar f_{U,s}^{\,*}\bigl(M_{(s),U}\bigr),
\]
con las notaciones de 2.3.1. Mostraremos que los \(\Proj_0(M)_{U,s}\) se pegan cuando varían \(U\) y \(s\), definiendo así un \(\mathcal O_X\)-módulo denotado por \(\Proj_0(M)\). Siendo evidente el pegado para \(U\) variable y \(s\) fijo, se reduce (cf. EGA II 2.5.2) a mostrar que, si \(t\) es una sección de \(B\) sobre \(U\), homogénea de grado \(e>o\), entonces la restricción de \(\Proj_0(M)_{U,s}\) a \(D_+(U,st)\) es canónicamente isomorfa a \(\Proj_0(M)_{U,st}\). Consideremos:
\[
\begin{tikzcd}[column sep=large,row sep=large]
D_+(st) \arrow[r,hook,"i"] \arrow[dr,"f_{U,st}"'] &
D_+(s) \arrow[d,"f_{U,s}"]\\
& U
\end{tikzcd}
\]
se trata de mostrar que los módulos
\[
  i^*\circ\bar f_{U,s}^{\,*}\bigl(M_{(s),U}\bigr)
  \quad\text{y}\quad
  \bar f_{U,st}^{\,*}\bigl(M_{(st),U}\bigr)
\]
son canónicamente isomorfos. Ahora bien, se tienen por una parte los isomorfismos canónicos:
\[
\begin{aligned}
i^*\circ\bar f_{U,s}^{\,*}\bigl(M_{(s),U}\bigr)
&\simeq
i^*\Bigl(f_{U,s}^{-1}\bigl(M_{(s),U}\bigr)
 \otimes_{f_{U,s}^{-1}((B|_U)_{(s)})}\mathcal O_{D_+(s)}\Bigr)\\
&\simeq
i^{-1}\Bigl(\sisi{i_{U,s}^{-1}}{f_{U,s}^{-1}}\bigl(M_{(s),U}\bigr)
 \otimes_{f_{U,s}^{-1}((B|_U)_{(s)})}\mathcal O_{D_+(s)}\Bigr)
 \otimes_{i^{-1}(\mathcal O_{D_+(s)})}\mathcal O_{D_+(st)}\\
&\simeq
f_{U,st}^{-1}\bigl(M_{(s),U}\bigr)
 \otimes_{f_{U,st}^{-1}((B|_U)_{(s)})}\mathcal O_{D_+(st)} .
\end{aligned}
\]
Por otra parte:
\[
  \bar f_{U,st}^{\,*}\bigl(M_{(st),U}\bigr)
  \simeq
  f_{U,st}^{-1}\bigl(M_{(st),U}\bigr)
  \otimes_{f_{U,st}^{-1}((B|U)_{(st)})}\mathcal O_{D_+(st)} .
\]
Por último, se deduce de EGA II 2.2.2, traducido en términos de haces, que se tienen los isomorfismos canónicos:
\[
  M_{(st),U}\simeq
  \bigl(M_{(s),U}\bigr)_{t^d/s^e}
  \simeq
  M_{(s),U}\otimes_{(B|U)_{(s)}}(B|U)_{(st)},
\]
de donde resulta finalmente el isomorfismo buscado.

Es claro que, si \(M\) es cuasicoherente, el módulo \(\Proj_0(M)\) así construido es precisamente el definido en EGA II.

